\documentclass[11pt,a4paper]{article}
\usepackage[utf8]{inputenc}
\usepackage[english]{babel}
\usepackage[T1]{fontenc}
\usepackage{geometry}
\geometry{a4paper, margin=1in}
\usepackage{hyperref}
\usepackage{booktabs}
\usepackage{amsmath}
\usepackage{titlesec}
\usepackage{fancyhdr}

\titleformat{\section}{\large\bfseries}{\thesection.}{0.5em}{}
\titleformat{\subsection}{\normalsize\bfseries}{\thesubsection.}{0.5em}{}

\pagestyle{fancy}
\fancyhf{}
\rhead{\small TaaSim Casablanca --- Project Summary}
\lfoot{\small ENSA Al Hoceima}
\rfoot{\small Page \thepage}

\title{\textbf{TaaSim Casablanca \\ An Intelligent Data-Driven Platform for Urban Mobility \\ \Large Technical Report Summary (3-Page Limit)}}

\author{
  \textbf{Realised by:} Salma Said \& Mohamed Tamzirt \\
  \textbf{Supervised by:} Prof. Mohamed El Marouani \\
  \small National School of Applied Sciences --- Al Hoceima (ENSAH) \\
  \small Big Data Capstone Project
}
\date{June 2026}

\begin{document}

\maketitle

\section{Introduction \& Problem Statement}
Casablanca, Morocco's bustling economic hub, is home to over 4 million residents who rely heavily on a highly fragmented public transit network. Petits taxis, grands taxis, and informal bus routes operate independently with zero coordination, no digital booking options, and no shared schedules. Up until now, city planners and operators have lacked a unified data layer to connect supply with rider demand in real time.

We built \textbf{TaaSim} (Transport as a Service -- Simulation) to address this gap. Our core idea is to treat urban mobility as a pure data engineering challenge. By ingesting real-time vehicle GPS streams, coordinating citizen ride reservations on the fly, and feeding historical batch data into machine learning pipelines, TaaSim bridges the gap between riders and drivers. Specifically, the platform enables:
\begin{itemize}
    \item \textbf{Dynamic rider-to-vehicle matching} (under 5 seconds P95 matching latency).
    \item \textbf{Surge demand forecasting} per zone 30 minutes in advance.
    \item \textbf{A unified analytical dashboard} giving city planners direct insights into transit bottlenecks.
\end{itemize}

\section{Methodology}
TaaSim relies on a modern Kappa Architecture and geometric coordinate projections to drive Casablanca-specific analytics from historical trip datasets.

\subsection{A Unified Kappa Architecture}
Rather than maintaining separate codebases for batch and streaming layers (the classic Lambda approach), we chose a \textbf{Kappa Architecture}. Apache Kafka acts as our central event bus (system of record) with a 7-day retention period. Apache Flink serves as our primary streaming and stateful processing engine. For batch analytics and machine learning model training, Apache Spark draws data directly from the S3-compatible MinIO data lake, keeping our codebase clean, unified, and free of batch-stream logic drift.

\subsection{BBox Affine Mapping (Porto $\rightarrow$ Casablanca)}
Since real-world transit GPS datasets for Casablanca are not publicly available, we designed a relative-position affine transformation (ADR-01) to project Porto's taxi dataset (1.7M trips) onto Casablanca's actual geographic bounds:
\begin{align*}
rel\_lon &= \frac{lon - Porto.min\_lon}{Porto.max\_lon - Porto.min\_lon} \\
rel\_lat &= \frac{lat - Porto.min\_lat}{Porto.max\_lat - Porto.min\_lat} \\
Casa\_lon &= Casa.min\_lon + rel\_lon \times (Casa.max\_lon - Casa.min\_lon) \\
Casa\_lat &= Casa.min\_lat + rel\_lat \times (Casa.max\_lat - Casa.min\_lat)
\end{align*}
To ensure coordinate integrity, we clamp relative coordinates to $[0,1]$ and run data-quality gates at the ingestion phase to verify that less than 1\% of points fall outside boundaries.

\subsection{Real-Time Stateful Streaming (Flink)}
We deployed three cooperative Flink streaming jobs to handle live event processing:
\begin{itemize}
    \item \textbf{Job 1 (GPS Normalizer):} Ingests raw GPS pings from Kafka, applies watermarks (managing up to 3 minutes of out-of-order data), and anonymizes coordinates by snapping them to arrondissement centroids before saving them to Apache Cassandra.
    \item \textbf{Job 2 (Demand Aggregator):} Counts active vehicles and pending ride requests within 30-second tumbling windows, continuously updating the city-wide supply/demand ratio.
    \item \textbf{Job 3 (Trip Matcher):} Matches waiting riders with the nearest available same-zone taxi using a stateful RocksDB backend, dynamically falling back to adjacent zones within 5 seconds if none are nearby.
\end{itemize}

\subsection{Batch Analytics \& Predictive ML (Spark)}
\begin{itemize}
    \item \textbf{Optimized Spark ETL:} To prevent memory exhaustion (OOM) on a 1.94 GB raw CSV file, we replaced slow schema inference with explicit type definitions (\texttt{StructType}) and converted broadcast joins to local conditional Column expressions (\texttt{when().otherwise()}), completing the job in under 5 minutes.
    \item \textbf{Surge Forecasting:} Using Spark MLlib, we trained a \texttt{GBTRegressor} on combined temporal (hours/days), spatial (arrondissement metadata), and weather (Open-Meteo API) features, outperforming the naive 7-day-lag baseline by predicting demand surges 30 minutes in advance.
\end{itemize}

\section{Implementation \& Achievements}
TaaSim is built as a hardened, production-ready stack deployed via Docker:
\begin{enumerate}
    \item \textbf{Orchestration:} 12 containerized services running seamlessly with custom health checks and S3A jar file-system sharing between Spark and Flink.
    \item \textbf{ETL Optimization:} The PySpark Porto ETL job was optimized by removing shuffling and driver-side CSV reads. \textbf{Final run duration: 4.4 minutes (SLA target < 5 minutes met successfully)}.
    \item \textbf{API \& Security:} Implented a FastAPI backend secured via HTTPS (SSL/TLS certificates) and role-based JWT access tokens, supported by Kafka ACLs preventing unauthorized write access to internal topics.
    \item \textbf{Visualization:} Built interactive Grafana dashboard panels querying Cassandra directly to display real-time taxi locations, demand heatmaps, and batch KPI summaries.
\end{enumerate}

\section{Experimental Results \& SLA Verification}
Locust load testing and automated integration metrics verified that the platform is ready for metropolitan deployment:

\begin{table}[h!]
\centering
\caption{SLA Performance Verification}
\begin{tabular}{@{}llccr@{}}
\toprule
\textbf{SLA / Metric} & \textbf{Target} & \textbf{Measured Performance} & \textbf{Status} \\ \midrule
Trip matching latency & < 5 seconds (P95) & 0.42 seconds & Met \\
GPS position freshness & < 15 seconds & 1.2 seconds & Met \\
Demand zone aggregate frequency & Every 30 seconds & 30 seconds & Met \\
Forecast API response time & < 500 milliseconds (P95) & 35 milliseconds & Met \\
Spark ETL Porto (1.7M rows) & < 5 minutes & 4.4 minutes (265.5s) & Met \\ \bottomrule
\end{tabular}
\end{table}

\section{Conclusion}
TaaSim Casablanca establishes a robust Big Data framework for modernizing municipal transit. By fusing stateful stream processing (Flink) with high-throughput batch ETLs (Spark) and low-latency database queries (Cassandra), the platform satisfies all required non-functional SLAs. It stands ready to provide riders with fast dispatching and city planners with the insights needed to optimize urban mobility.

\end{document}
